\documentclass[10pt,a4paper]{amsart}
\usepackage{amsmath,amssymb,amsthm,graphicx}
\numberwithin{equation}{section}

\newtheorem{Theorem}{Theorem}[section]
\newtheorem*{Theorem*}{Theorem}
\newtheorem{Corollary}[Theorem]{Corollary}
\newtheorem{Lemma}[Theorem]{Lemma}
\newtheorem{Proposition}[Theorem]{Proposition}

\theoremstyle{definition}
\newtheorem{Definition}[Theorem]{Definition}
\newtheorem{Note}[Theorem]{Note}
\newtheorem{Example}[Theorem]{Example}
\newtheorem{Remark}[Theorem]{Remark}

\newcommand{\GL}{\operatorname{GL}}
\newcommand{\Hom}{\operatorname{Hom}}
\newcommand{\Mor}{\operatorname{Mor}}
\newcommand{\Map}{\operatorname{Map}}
\newcommand{\Ad}{\operatorname{Ad}}
\newcommand{\Aut}{\operatorname{Aut}}
\newcommand{\Isom}{\operatorname{Isom}}
\newcommand{\Out}{\operatorname{Out}}
\newcommand{\Inn}{\operatorname{Inn}}
\newcommand{\spec}{\operatorname{Spec}}
\newcommand{\mon}{\operatorname{mon}}
\newcommand{\Hol}{\operatorname{Hol}}
\newcommand{\Lie}{\operatorname{Lie}}
\newcommand{\curv}{\operatorname{curv}}
\newcommand{\pr}{\operatorname{pr}}
\newcommand{\ave}{\operatorname{ave}}
\newcommand{\bG}{\mathbb{G}}
\newcommand{\bH}{\mathbb{H}}
\newcommand{\bD}{\mathbb{D}}
\newcommand{\bC}{\mathbb{C}}
\newcommand{\bA}{\mathbb{A}}
\newcommand{\bB}{\mathbb{B}}
\newcommand{\bR}{\mathbb{R}}
\newcommand{\bZ}{\mathbb{Z}}
\DeclareMathOperator*{\bigfusion}{\text{\rlap{$\bigoplus$}$\bigotimes$}}
\usepackage[bookmarks=false]{hyperref}
\title{Quasi-Hamiltonian structures on equivariant surface representation spaces and their twisted fixed-point spaces}
\author{Keito Takegoshi}
\date{doi:10.3842/SIGMA.2026.088; CC BY 4.0}
\begin{document}
\maketitle
\noindent\textit{Allowed standard prerequisites (editorial).}\begin{itemize}
\item Alekseev-Malkin-Meinrenken quasi-Hamiltonian/loop-group correspondence. AMM, author cites Theorems 8.3 and 9.1 with source-local construction retained
\item Sobolev gauge theory on compact surfaces with boundary. Author section 5 and bibliography
\end{itemize}\medskip
\section{Finite group actions on quasi-Hamiltonian spaces}\label{2Fin_grp_act_q-Ham_sp}
In this section, we give the definition of a~finite group $\Gamma$-compatible $q$-Hamiltonian $G$-space and present a~basic example that is an analogue of the double.

\subsection{Compatibilities of bitorsors}

Let $G$ be a~compact Lie group with Lie algebra $\mathfrak{g}$ and let $\Gamma$ be a~finite group. We fix an action $\Gamma \to \Aut(G)$ of $\Gamma$ on $G$.
Also we give a~$G$-bitorsor $\bG$ equipped with a~left $\Gamma$-action where a~bitorsor $\bG$ is a~manifold endowed with left and right $G$-actions that are both simply transitive and commute with each other.

\begin{Definition} A~$\Gamma$-compatible $G$-bitorsor $\bG$ is a~bitorsor $\bG$ equipped with a~$\Gamma$-action such that $
 \varphi \cdot (g x) = (\varphi \cdot g)(\varphi \cdot x)$, $
 \varphi \cdot (x g) = (\varphi \cdot x)(\varphi \cdot g)$, ($
 \varphi \in \Gamma$, $ g \in G$, $ x \in \bG$).
\end{Definition}

\begin{Example}\label{exam_bitorsor}
 We define the space $\bG := \Map(\bZ/ m \bZ, G)$.
 Fixing an action $\varphi \colon \bZ/m\bZ \to \Aut(G)$, we define a~$\bZ / m \bZ $-action on $\bG$ by
$(l \cdot g)(k) := \varphi(l) g (k - l)$, $ l, k \in \bZ/ m \bZ$, $ g \in \bG$.
 Then the space~$\bG$ is a~$\bZ/ m \bZ$-compatible $\Map(\bZ/ m \bZ, G)$-bitorsor.
\end{Example}

\begin{Remark}\label{compati_bitorsor_and_semi_product}
 Let $\bG$ be a~$\Gamma$-compatible $G$-bitorsor.
 Define left and right actions of~$G \rtimes \Gamma$ on~$\bG \times \Gamma$ by $(g,\varphi)(x,\psi) := (g(\varphi \cdot x), \varphi \psi)$, $(x, \psi)(g,\varphi) := (x(\psi \cdot g), \psi \varphi)$,
 where $(g,\varphi)\in G \rtimes \Gamma$ and~${(x, \psi) \in \bG \times \Gamma}$.
 Then these actions make $\bG \times \Gamma$ into a~$G \rtimes \Gamma$-bitorsor.
 We denote this bitorsor by $\bG \rtimes \Gamma$.
\end{Remark}

We define the inverse $\bG^{-1}$ of a~$G$-bitorsor $\bG$ by the space $\bG^{-1} := \bigl\{x^{-1}\mid x \in \bG\big\}$, which is~a~copy of $\bG$, equipped with the following $G$-actions,
\smash{$ g x^{-1} := \bigl(x g^{-1}\bigr)^{-1}$}, \smash{$ x^{-1} g := \bigl(g^{-1} x\bigr)^{-1}$}
for~${g \in G}$ and $x \in \bG$.

Moreover, we define the product $\bG_1\cdot \bG_2$ of two $G$-bitorsors $\bG_1$ and $\bG_2$ as the quotient of~$\bG_1 \times \bG_2$ by the following action,
$g (x_1,x_2) := \bigl(x_1 g^{-1}, g x_2\bigr)$
for $g \in G$ and $(x_1,x_2) \in \bG_1 \times \bG_2$.
We denote an element $[x_1,x_2] \in \bG_1 \cdot \bG_2$ by $x_1 \cdot x_2$.

\begin{Proposition}
Let $\bG$, $\bG_1$, $\bG_2$ be $\Gamma$-compatible $G$-bitorsors. Suppose that the fixed point sets~$\bG^{\Gamma}$,~$\bG_1^{\Gamma}$,~$\bG_2^{\Gamma}$ are non-empty.
Then we have the following properties:
\begin{enumerate}\itemsep=0pt
 \item [$(1)$] The fixed point set $\bG^{\Gamma}$ becomes a~$G^{\Gamma}$-bitorsor under the restriction of the action of $G$ on~$\bG$.
 \item [$(2)$] The inverse $\bG^{-1}$ of the bitorsor $\bG$ and the product $\bG_1 \cdot \bG_2$ of $\bG_1$ and $\bG_2$ become $\Gamma$-compatible $G$-bitorsors.
 \item [$(3)$] There exists an isomorphism of $G^{\Gamma}$-bitorsor $\bG_1^{\Gamma}\cdot \bG_2^{\Gamma} \simeq (\bG_1 \cdot \bG_2)^{\Gamma}$.
\end{enumerate}
\end{Proposition}

\begin{proof}
 Statement (1) is easy to verify.
 We define actions of $\Gamma$ on \smash{$\bG_1^{-1}$} and $\bG_1\cdot \bG_2$ by \smash{$\varphi x^{-1} := (\varphi x)^{-1}$}, $\varphi (x_1 \cdot x_2) := (\varphi x_1) \cdot (\varphi x_2)$,
 for $\varphi \in \Gamma$ and $x_j \in \bG_j$.
 These actions satisfy the definition of $\Gamma$-compatibility.

 To show statement (3), we see that the canonical map $ \Phi \colon \bG_1^{\Gamma} \cdot \bG_2^{\Gamma} \to (\bG_1 \cdot \bG_2)^{\Gamma}$, $ x_1 \cdot x_2 \mapsto x_1 \cdot x_2$
 is equivariant with respect to the left and right $G^{\Gamma}$-actions.
 For any $g \in G^{\Gamma}$, $x_1 \cdot x_2 \in \bG_1^{\Gamma} \cdot \bG_2^{\Gamma}$,
 \begin{align*}
 \Phi(g (x_1 \cdot x_2)) &= \Phi(g x_1 \cdot x_2) = gx_1 \cdot x_2 = g (x_1 \cdot x_2) = g \Phi(x_1 \cdot x_2).
 \end{align*}
 The equivariance of $\Phi$ with respect to the right $G^{\Gamma}$-action is verified in the same way as above.
 Hence, by the simple transitivity of the left and right actions on bitorsors, $\Phi$ is an isomorphism.\looseness=1
\end{proof}

\subsection[Compatibilities of q-Hamiltonian G-spaces]{Compatibilities of $\boldsymbol{q}$-Hamiltonian $\boldsymbol{G}$-spaces}

First, we recall the definition of twisted quasi-Hamiltonian $G$-spaces~\cite{B-Y2015,Meinrenken2017}.
Throughout this paper, we fix an $\Ad$-invariant positive definite inner product $(\ ,\ )$ on $\mathfrak{g}$ and consider only bitorsors whose corresponding elements in $\Out(G)$ preserve this inner product.
Also we assume that the $\Gamma$-action on $\mathfrak{g}$ induced by the action on $G$ preserves the inner product on $\mathfrak{g}$.

We define the left Maurer--Cartan form $\theta$ and the right Maurer--Cartan form~$\overline{\theta}$ on a~$G$-bitorsor~$\bG$ by
$\theta_x(X) := x^{-1}X$, $ \overline{\theta}_x(X) := Xx^{-1}$, $ x \in \bG$, $ X \in T_x\bG$,
where $x^{-1}X \in \mathfrak{g}$ denotes the derivative of the map $x^{-1} \colon \bG \to G$, $ y \mapsto x^{-1}y$ for $X$.

\begin{Definition} A~(twisted) quasi-Hamiltonian $G$-space $(M, \omega, \mu)$ is a~$G$-manifold $M$ together with an invariant $2$-form $\omega$ and a~$G$-equivariant map $\mu \colon M \to \bG$ to a~$G$-bitorsor $\bG$ such that
\begin{itemize}\itemsep=0pt\setlength{\leftskip}{0.30cm}
 \item[(QH1)] ${\rm d}\omega = \mu^* \frac{1}{12}(\theta,[\theta,\theta])$.
 \item[(QH2)] $\iota \bigl(\xi ^{\#}\bigr)\omega = \frac{1}{2} \mu^* \bigl(\theta + \overline{\theta},\xi\bigr) $, ($\xi \in \mathfrak{g})$, $\xi^{\#}|_p = \frac{\rm d}{{\rm d} t}\bigl({\rm e}^{-t\xi}\cdot p\bigr)|_{t=0}$.
 \item[(QH3)] $\ker \omega_p \cap \ker d\mu_p =\{0\}$, ($p \in M$),
 \end{itemize}
where $\theta$ (resp.\ $\overline{\theta}$) denotes the left (resp.\ right) Maurer--Cartan form on $\bG$.
\end{Definition}

We note that the restriction of the inner product $(\ ,\ )$ to the Lie algebra $\mathfrak{g}^{\Gamma}$ of the $\Gamma$-fixed point Lie subgroup $G^{\Gamma} := \{g \in G\mid \varphi \cdot g = g \text{ for all } \varphi \in \Gamma\}$ is also positive definite and invariant.
Now we give the definition of a~$\Gamma$-compatible quasi-Hamiltonian $G$-space.

\begin{Definition}\label{def_compati_q-Ham}
 A~$\Gamma$-compatible quasi-Hamiltonian $G$-space $(M, \omega, \mu)$ is a~twisted quasi-Hamil\-tonian $G$-space equipped with a~$\Gamma$-action such that
\begin{enumerate}\itemsep=0pt
 \item[$(1)$] $\varphi \cdot (g \cdot p) = (\varphi \cdot g) \cdot (\varphi \cdot p)$, ($\varphi \in \Gamma$, $g \in G$, $ p \in M$).
 \item[$(2)$] The $q$-Hamiltonian $2$-form $\omega$ is $\Gamma$-invariant.
 \item[$(3)$] The moment map $\mu \colon M \to \bG$ is $\Gamma$-equivariant.
 \end{enumerate}
We say that any two $\Gamma$-compatible quasi-Hamiltonian $G$-spaces are isomorphic if there exists a~$\Gamma$-equivariant isomorphism of twisted quasi-Hamiltonian $G$-spaces.
\end{Definition}

From a~$\Gamma$-compatible quasi-Hamiltonian $G$-space, one can naturally construct a~twisted quasi-Hamiltonian $G \rtimes \Gamma$-space.
\begin{Remark}\label{cpti_q-Ham_and_twisted_q-Ham}
 Let $(M,\omega,\mu)$ be a~$\Gamma$-compatible quasi-Hamiltonian $G$-space whose moment map takes values in a~$G$-bitorsor $\bG$.
 The actions of $G$ and $\Gamma$ naturally induce an action of $G \rtimes \Gamma$ on~$M$ given by $(g,\varphi) \cdot p := g \cdot (\varphi \cdot p) $, ($(g,\varphi) \in G \rtimes \Gamma$, $ p \in M$),
 and a~moment map $\tilde{\mu} := M \to \bG \rtimes \Gamma, p \mapsto (\mu(p),1_{\Gamma})$,
 where $\bG \rtimes \Gamma$ is the bitorsor introduced in~Remark~\ref{compati_bitorsor_and_semi_product}.
 Then we obtain a~twisted quasi-Hamiltonian $G \rtimes \Gamma$-space $(M,\omega,\tilde{\mu})$.
\end{Remark}

\subsection[The double D(G\_1,G\_2)]{The double $\boldsymbol{\text{D}(\bG_1,\bG_2)}$}\label{the_double}

We give a~basic example of a~$\Gamma$-compatible quasi-Hamiltonian space.
Let $\bG_1$ and $\bG_2$ be $\Gamma$-compatible $G$-bitorsors.
We define a~$G \times G$-action on $D := \bG_1 \times \bG_2$ by
\[ (g_1, g_2) \cdot (a,b) := \bigl(g_1 a g_2^{-1}, g_2 b g_1^{-1}\bigr), \qquad (g_1,g_2) \in G \times G,\qquad (a,b) \in D.\]
Also we define a~$2$-form $\omega$ and a~map $\mu$ by
\begin{gather*}
 \omega := -\frac{1}{2}\bigl(a^* \theta, b^* \overline{\theta}\bigr)-\frac{1}{2}\bigl(a^*\overline{\theta},b^*\theta\bigr),\\
 \mu \colon \ D \to \bG_1\cdot \bG_2 \times \bG_1^{-1} \cdot \bG_2^{-1},\qquad
(a,b) \mapsto \bigl(ab, a^{-1}b^{-1}\bigr).
\end{gather*}
The tuple $\textbf{D}(\bG_1,\bG_2) := (D,\omega, \mu)$ is a~twisted quasi-Hamiltonian $(G\times G)$-space.
Moreover, under the diagonal action of $\Gamma$ on $D$ and $\bG_1\bG_2 \times \bG_1^{-1}\bG_2^{-1}$, the space $\text{D}(\bG_1,\bG_2)$ carries a~$\Gamma$-compatible structure.

\begin{Remark}\label{the double with different coordinates}
 Let $\bC$ and $\bH$ be $\Gamma$-compatible $G$-bitorsors. We give the double $\textbf{D}(\bC,\bH)$ with different coordinates.
 For any $(C,h) \in \textbf{D}(\bC,\bH)$ and $(g,k) \in G \times G$, we define an action of $G \times G$ on $\textbf{D}(\bC,\bH)$ by $(g,k) \cdot (C,h) := \bigl(kCg^{-1},khk^{-1}\bigr)$.
 We also define a~quasi-Hamiltonian $2$-form by
 \[\omega := \frac{1}{2}\bigl(C^*\overline{\theta},\Ad_{h}C^*\overline{\theta}\bigr)+\frac{1}{2}\bigl(C^*\overline{\theta},h^*\theta+h^*\overline{\theta}\bigr),\]
 and a~moment map by
 \[
 \mu \colon\ \textbf{D}(\bC,\bH) \to \bC^{-1} \bH \bC \times \bH^{-1},\qquad (C,h) \mapsto \bigl(C^{-1}hC, h^{-1}\bigr).
 \]
 Then the twisted quasi-Hamiltonian $(G \times G)$-space $\textbf{D}(\bC,\bH)$ together with the natural $\Gamma$-action is also a~$\Gamma$-compatible quasi-Hamiltonian $(G \times G)$-space and is isomorphic to the double of $\bC^{-1}$ and~$\bH \bC$ with standard coordinates.
\end{Remark}

\section{Basic properties}\label{3Basic_properties}

Twisted quasi-Hamiltonian spaces have been investigated by several authors~\cite{B-Y2015,Knop,Meinrenken2017}.
Although a~$\Gamma$-compatible quasi-Hamiltonian $G$-space naturally gives rise to a~twisted quasi-Hamil\-tonian $(G \rtimes \Gamma)$-space (see Remark~\ref{cpti_q-Ham_and_twisted_q-Ham}), we study some basic properties of $\Gamma$-compatible quasi-Hamiltonian $G$-spaces, since this is the natural framework for the applications to twisted moduli spaces.

\subsection[The set of fixed points of (M, omega, mu)]{The set of fixed points of $\boldsymbol{(M, \omega, \mu)}$}

We fix an action $\Gamma \to \Aut(G)$.
Let $(M,\omega, \mu)$ be a~$\Gamma$-compatible $q$-Hamiltonian $G$-space with the moment map $\mu \colon M \to \bG$.
Let $M^{\Gamma}$ be the fixed point set of $M$ for the action of $\Gamma$ and let $\bG^{\Gamma}$ be the fixed point set of $\bG$.
We denote the identity component of $G^{\Gamma}$ by $H$.
Fix a~connected component $N$ of $M^{\Gamma}$, and we can then take the connected component $\bH$ of $\bG$ so that $\mu(N) \subset \bH$.
Then $\bH$ becomes an~$H$-bitorsor and the following theorem holds.

\begin{Theorem}\label{fixed_points_q-Ham}
 The tuple $(N, \omega|_N, \mu|_N)$ is a~quasi-Hamiltonian $H$-space.
\end{Theorem}

\begin{proof}Let $\theta^{\Gamma}$ (resp.\ \smash{$\overline{\theta}^{\Gamma}$}) be the left (resp.\ right) Maurer--Cartan form on $\bH$.
 We denote the inclusions by $\iota_N \colon N \to M$ and $\iota_{\bH} \colon \bH \to \bG $.
Let the Cartan $3$-forms on $\bG$ and $\bH$ be~\smash{${\chi := \frac{1}{12}(\theta,[\theta,\theta])_{\mathfrak{g}}}$}
 and \smash{$\chi^{\Gamma} := \frac{1}{12}\bigl(\theta^{\Gamma},\bigl[\theta^{\Gamma},\theta^{\Gamma}\bigr]\bigr)_{\mathfrak{g}^{\Gamma}}$},
 respectively.
 Since \smash{$\chi^{\Gamma} = \chi |_{\bH}$}, we obtain the following $
 {\rm d}\omega|_N
 = \bigl(\mu^{*}\chi\bigr)|_{N}
 = (\mu|_N)^* \chi^{\Gamma}$.

Next, we will show the moment map condition.
 For any $\xi \in \mathfrak{g}^{\Gamma}$,
 \begin{align*}
 \iota\bigl(\xi^{\#}\bigr)\omega|_N &= \bigl(\iota\bigl(\xi^{\#}\bigr)\omega\bigr)|_{N} = \left.\left( \frac{1}{2} \mu^*\bigl(\theta + \overline{\theta},\xi \bigr) _{\mathfrak{g}}\right)\right| _{N} = \frac{1}{2}{\mu|_N}^*\bigl(\theta^{\Gamma}+\overline{\theta}^{\Gamma},\xi\bigr)_{\mathfrak{g}^{\Gamma}}.
 \end{align*}

 The degeneracy condition of $\omega|_N$ is shown by using the averaging over $\Gamma$.
 For any $x \in N$ and~${X \in \ker{\omega|_N}_x \cap \ker {\rm d}_{x}\mu|_N}$,
${\rm d}\mu \circ {\rm d}\iota_{N}(X) = {\rm d}\iota_{\bH} \circ {\rm d} \mu|_N(X)= 0$.
 This leads to ${\rm d}\iota_{N}(X) \in \ker {\rm d}_{x}\mu$.
 For any $Y \in T_xM$, the following identity holds from the~$\Gamma$-invariance of~$\omega$
 \begin{align*}
 \omega ({\rm d}\iota_{N}(X), Y) &= \frac{1}{|\Gamma|}\sum_{\varphi \in \Gamma} \varphi^*\omega({\rm d}\iota_{N}(X),Y) = \omega\left({\rm d}\iota_{N}(X), \frac{1}{|\Gamma|}\sum \varphi_*Y\right) =0.
 \end{align*}
Therefore, we obtain ${\rm d}\iota_{N}(X) \in \ker \omega_x \cap \ker {\rm d}_x \mu = \{0\}$.
So, we conclude
\[
\ker\omega|_N \cap \ker {\rm d}\mu|_N =\{0\}.\tag*{\qed}
\]\renewcommand{\qed}{}
\end{proof}

\subsection[Fusion products and G-compatibilities]{Fusion products and $\boldsymbol{\Gamma}$-compatibilities}

Let $G$, $H_1$, and $H_2$ be compact Lie groups.
Let $\bG$, $\bH_1$, and $\bH_2$ be bitorsors for $G$, $H_1$, and $H_2$ respectively.
We fix actions of a~finite group $\Gamma$ on $G$ and $H_j$.
Let $(M_j, \omega_j, \mu_j)$ be a~$\Gamma$-compatible quasi-Hamiltonian $(G\times H_j)$-space for $j=1,2$.
We denote the moment map $\mu_j \colon M_j \to \bG_j \times \bH_j$ by $\mu_j =\bigl(\mu^{G}_j,\mu^{H}_j\bigr)$.

\begin{Proposition}
 The fusion product $M_1 \circledast M_2$ carries a~$\Gamma$-compatible quasi-Hamiltonian $G \times H_1 \times H_2$-structure for the diagonal action of $\Gamma$ on $M_1 \circledast M_2$, $G_1 \times H_1 \times H_2$ and $\bG_1 \bG_2 \times \bH_1 \times \bH_2$.
\end{Proposition}

\begin{proof}
 We show the quasi-Hamiltonian $2$-form on $M_1 \circledast M_2$, $\omega = \omega_1 + \omega_2 - \frac{1}{2}\bigl(\bigl(\mu_1^{G}\bigr)^*\theta, \bigl(\mu_2^{G}\bigr)^* \overline{\theta}\bigr)$
 is $\Gamma$-invariant. Since $\bG_j$ are $\Gamma$-compatible and the inner product $(\ ,\ )$ on $\mathfrak{g}$ is invariant under the action of $\Gamma$, for any $\varphi \in \Gamma$,
 \begin{align*}
 \varphi^*\omega &= \varphi ^* \omega_1 + \varphi ^* \omega_2 -\frac{1}{2}\varphi^*\bigl(\bigl(\mu_1^G\bigr)^*\theta,\bigl(\mu_2^G\bigr)^*\overline{\theta}\bigr) = \omega_1 + \omega_2 -\frac{1}{2}\bigl(\varphi \cdot \bigl(\mu_1^G\bigr)^*\theta, \varphi \cdot \bigl(\mu_2^G\bigr)^* \overline{\theta}\bigr) \\
 &= \omega_1 + \omega_2 -\frac{1}{2}\bigl(\bigl(\mu_1^G\bigr)^*\theta, \bigl(\mu_2^G\bigr)^* \overline{\theta}\bigr) = \omega.
 \end{align*}
 The conditions $(1)$ and $(3)$ in Definition~\ref{def_compati_q-Ham} are easy to verify in this case.
\end{proof}

Moreover, suppose that $M_j^{\Gamma}$ are non-empty connected sets and $\bG_j^{\Gamma}$ and $\bH_j^{\Gamma}$ are non-empty.
One can easily prove the following proposition.

\begin{Proposition}\label{fusion_and_fixedpoints}
 There exists a~canonical twisted quasi-Hamiltonian isomorphism $M_1^{\Gamma} \circledast M_2^{\Gamma} \simeq (M_1 \circledast M_2)^{\Gamma}$.
\end{Proposition}

\begin{Remark}
 In Proposition~\ref{fusion_and_fixedpoints}, if $M_j^{\Gamma}$ are not connected, we can obtain a~similar result as above by taking connected components of $M_1^{\Gamma}$ and $M_2^{\Gamma}$.
\end{Remark}

\begin{Example}[fused double $\bD(\bG_1,\bG_2)$]
 Let $\bG_1$ and $\bG_2$ be $\Gamma$-compatible $G$-bitorsors.
 The internal fusion of the double $\textbf{D}(\bG_1,\bG_2)$ (see Example~\ref{the_double}) is denoted by $\bD(\bG_1,\bG_2)$, and is called the fused double.
 The space $\bD(\bG_1,\bG_2)$ carries a~$\Gamma$-compatible structure.
 Moreover, $\bG_1^{\Gamma}$ and $\bG_2^{\Gamma}$ are submanifolds of $\bG_1$ and $\bG_2$, respectively.
 We obtain an isomorphism of twisted quasi-Hamiltonian $G^{\Gamma}$-space
$\bD(\bG_1,\bG_2)^{\Gamma} \simeq \bD\bigl(\bG_1^{\Gamma},\bG_2^{\Gamma}\bigr)$.
\end{Example}

\section{Cyclic group actions on loop groups}\label{4Cyclic_grp_act}

The results of this section provide the foundation for the constructions in Section~\ref{5Tw_conn_and_fin_grp_act}.
We introduce the notion of Hamiltonian $LG$-spaces compatible with an action of a~cyclic group $\Gamma := \bZ/m\bZ$, examine their relationship with Hamiltonian $L^{(\kappa)}G$-spaces, and discuss the equivalence theorem in the $\Gamma$-compatible setting.
The constructions and computations involving the canonical $2$\nobreakdash-form on~$LG$ follow the same arguments as those in the (twisted) loop group setting developed in~\cite{AMM,LYM_twisted_LG}.

\subsection{Twisted loop groups}\label{Twisted loop groups}
Let $G$ be a~compact connected simply connected Lie group
with Lie algebra $\mathfrak{g}$. We fix an $\Ad$-invariant inner product $(\ ,\ )$ on $\mathfrak{g}$
and a~Lie group automorphism $\kappa \in \Aut(G)$ such that there exists a~positive integer $m$ such that $\kappa^m=1$.
We suppose that $\kappa$ preserves the inner product~$(\ ,\ )$ on $\mathfrak{g}$.
Let $S^{1}_m$ denote $\bR/ m \bZ$.
We define the loop group $LG$ as a~space of maps
\[LG := \Map\bigl(S^1_m,G\bigr)\]
of a~fixed Sobolev class $\lambda > 1/2$.
Also we define the $\kappa$-twisted loop group $L^{(\kappa)}G$ as a~subgroup~of~$LG$
\[L^{(\kappa)}G := \bigl\{g \in LG \mid \kappa g (t) = g(t+1) \text{ for all } t \in S^1_m\bigr\}.\]
We define the spaces $L^{(\kappa)}\mathfrak{g}$ and $L^{(\kappa)}\mathfrak{g}^*$ as subspaces of the space of maps $L\mathfrak{g}:= \Map\bigl(S^1_m, \mathfrak{g}\bigr)$ of~Sobolev class $\lambda$ and the space of $\mathfrak{g}$-valued $1$-forms $L\mathfrak{g}^* := \Omega^1\bigl(S^1_m,\mathfrak{g}\bigr)$ of Sobolev class $\lambda-1$,
\begin{gather*}
 L^{(\kappa)}\mathfrak{g} := \bigl\{\xi \in L\mathfrak{g}\mid \kappa \xi (t) = \xi (t+1) \text{ for all }t \in S^1_m \bigr\}, \\
 L^{(\kappa)}\mathfrak{g}^{*} := \bigl\{A \in L\mathfrak{g}^*\mid \kappa A_t = A_{t+1} \text{ for all }t \in S^1_m\bigr\}.
\end{gather*}
The gauge transformation on $L^{(\kappa)}\mathfrak{g}^*$ is defined by $g \cdot A := gAg^{-1} + {\rm d}gg^{-1}$, $ g \in L^{(\kappa)}G$, $ A \in L^{(\kappa)}\mathfrak{g}^*$.
There exists a~canonical pairing $L^{(\kappa)}\mathfrak{g}^* \times L^{(\kappa)}\mathfrak{g} \to \bR$, \smash{$ (A, \xi) \mapsto \int_{0}^{1}(A,\xi)$}.
We define the holonomy~\smash{$\Hol^{b}_t(A)$} of $A$ starting at $b \in S^{1}_m$ by the unique solution of the following differential equation:
\begin{align*}
 \Hol^{b}_t(A)^{-1} \frac{\partial \Hol^{b}_t(A)}{\partial t} &= -A,\qquad
 \Hol^{b}_b(A) = 1_G.
\end{align*}
Then, the holonomy and the gauge action satisfy the following equation:
\[
\Hol^{b}_t(g\cdot A) = g(b) \Hol ^{b}_t(A) g(t)^{-1}, \qquad g \in L^{(\kappa)}G,\qquad A \in L^{(\kappa)}\mathfrak{g}^*.
\]
Therefore, we have the holonomy map \smash{$\Hol^{(\kappa)} \colon L^{(\kappa)}\mathfrak{g}^{*} \to G^{(\kappa)}$}, \smash{$ A \mapsto \Hol^{0}_1(A)$},
where $G^{(\kappa)} = G$ is a~$G$-bitorsor with the following left and right actions
$g\cdot x := gx$, $x \cdot g := x (\kappa g)$, \big(\smash{$x \in G^{(\kappa)}$}, $ g \in G$\big).

Let $\theta$ (resp.\ $\overline{\theta}$) be the left (resp.\ right) Maurer--Cartan form on $G$.
We define the $2$-form on~$L^{(\kappa)}\mathfrak{g}^*$ by
\[
 \varpi ^{(\kappa)} := \frac{1}{2} \int_{0}^{1}{\rm d}s \left({\Hol^{0}_s}^* \overline{\theta}, \frac{\partial}{\partial s}{\Hol^{0}_s}^* \overline{\theta}\right).
 \]

The following proposition is analogous to~\cite[Proposition 8.1]{AMM}.

\begin{Proposition}We have
\begin{gather*}
 \varpi^{(\kappa)} \text{ is } L^{(\kappa)}G \text{-invariant},\qquad
 {\rm d} \varpi^{(\kappa)} = - {\Hol^{(\kappa)}}^* \frac{1}{12} (\theta,[\theta,\theta]),\\
 \iota\bigl(\xi^{\#}\bigr)\varpi^{(\kappa)} = - {\rm d} \int_{0}^{1}(A,\xi) - \frac{1}{2} {\Hol^{(\kappa)}}^*\bigl(\theta+\overline{\theta},\xi\bigr), \qquad \xi \in L^{(\kappa)}\mathfrak{g}.
\end{gather*}
\end{Proposition}

\begin{proof}
This is shown in the same way as in Proposition~\ref{formula on loop group}.
\end{proof}

\subsection{Cyclic group actions}

In this subsection, let $\Gamma$ be a~cyclic group $\bZ / m \bZ$, and let $S^{1}$ denote $\bR / m \bZ$.
We consider a~$\Gamma$-action on $G$, i.e., a~group homomorphism $\Gamma \to \Aut(G)$. Also a~natural $\Gamma$-action on $S^{1}$ is defined by
$ t \cdot l := t -l $ for $t \in S^{1}$ and $l \in \Gamma$.
Then we can define a~$\Gamma$-action on the loop group~$LG$ by
\[
l \cdot g (t) := \kappa^{l} g(t \cdot l) ,\qquad g \in LG,\qquad t \in S^{1},\qquad l \in \Gamma.
\]
$\Gamma$-actions on $L\mathfrak{g}$ and $L\mathfrak{g}^{*}$ are defined in the same way.
Let $\beta$ be a~finite subset $0\cdot \Gamma = \{0,1,\dots,\allowbreak {m-1}\} \subset S^{1}$, and we denote $G^{\beta} := \Map(\beta,G)$.
Also we denote Maurer--Cartan forms on $G^{\beta}$ by~$\theta$ and~$\overline{\theta}$.
Similar to the twisted case described above, gauge transformations, holonomies and natural pairing are defined.
Furthermore, we define the $2$-form on $L\mathfrak{g}^{*}$ by
\[
\varpi := \frac{1}{2} \sum_{i=1}^{m} \int_{i}^{i+1}{\rm d}s \left({\Hol^{i}_s}^{*}\overline{\theta},\frac{\partial}{\partial s}{\Hol^{i}_s}^{*}\overline{\theta}\right).
\]
Also the holonomy map $\Hol \colon L\mathfrak{g}^{*} \to \bH := G^{\beta}$ is defined by
\smash{$\Hol(A) := \bigl(\Hol^{i}_{i+1}(A)\bigr)^{m}_{i=1}$}, $\bigl(A \in L\mathfrak{g}^{*}\bigr)$,
where $\bH$ is a~$\Gamma$-compatible $G^{\beta}$-bitorsor (Example~\ref{exam_bitorsor}).

In this case, the following proposition holds in the same way as in~\cite[Proposition 8.1]{AMM}.

\begin{Proposition}\label{formula on loop group}We have
 \begin{gather*}
 \varpi\ \text{is}\ LG\text{-invariant}, \qquad
 {\rm d} \varpi = - \Hol ^{*}\frac{1}{12}(\theta,[\theta,\theta]), \\
 \iota\bigl(\xi^{\#}\bigr)\varpi_A = -{\rm d} \oint (A,\xi) - \frac{1}{2}\Hol^*\bigl(\theta + \overline{\theta}, \xi|_\beta\bigr), \qquad \xi \in L\mathfrak{g},\qquad A \in L\mathfrak{g}^*,
 \end{gather*}
 where $\xi|_{\beta} = (\xi(i))_{i \in \beta} \in \mathfrak{g}^{\beta}$.
\end{Proposition}

To prove it, we use the following lemma.

\begin{Lemma}\label{lemma of calc of Hol} We have
 \begin{gather*}
 \iota\bigl(\xi^{\#}\bigr){\Hol^{i}_s}^{*}\overline{\theta} = \Ad\bigl(\Hol^{i}_s\bigr)\xi(s) - \xi(i), \qquad \xi \in L\mathfrak{g}. \\
 \iota(\eta) {\Hol^{i}_s}^*\overline{\theta} = - \int_{i}^{s}{\rm d}u \Ad \bigl(\Hol^{i}_u\bigr)\eta , \qquad \eta \in T_A L\mathfrak{g}^{*}=L\mathfrak{g}^*.
 \end{gather*}
\end{Lemma}


\subsection[Gamma-compatibilities of Hamiltonian LG-spaces]{$\boldsymbol{\Gamma}$-compatibilities of Hamiltonian $\boldsymbol{LG}$-spaces}

In this subsection, we state the definition of a~$\Gamma$-compatibility on Hamiltonian $LG$-spaces.
This is the usual Hamiltonian space version of a~$\Gamma$-compatibility.

First, we recall the definition of a~Hamiltonian $LG$-space~\cite{AMM}.

\begin{Definition} A~Hamiltonian $LG$-space is a~Banach-manifold $N$ together with an $LG$-action, an invariant $2$-form \smash{$\sigma \in \Omega^2(N)^{LG}$}, and an equivariant map \smash{$\Phi \in C^{\infty}\bigl(N, L\mathfrak{g}^*\bigr)^{LG}$} such that
 \begin{itemize}\itemsep=0pt
 \item [(1)] The $2$-form $\sigma$ is closed.
 \item [(2)] The map $\Phi$ is a~moment map for the $LG$-action \smash{$\iota\bigl(\xi^{\#}\bigr)\sigma= -{\rm d} \oint_{S^1}(m, \xi)$}.
 \item [(3)] The $2$-form $\sigma$ is weakly non-degenerate, that is, the induced map \smash{$\sigma^{\flat}_x \colon T_xN \to T^{*}_xN$} is injective.
 \end{itemize}
\end{Definition}

\begin{Definition}
Let $(N,\sigma,\Phi)$ be a~Hamiltonian $LG$-space and we fix a~$\Gamma$-action on $N$.
 The tuple $(N,\sigma,\Phi)$ together with the $\Gamma$-action is called a~$\Gamma$-compatible Hamiltonian $LG$-space if it satisfies the following properties:
\begin{itemize}\itemsep=0pt
 \item[(1)] $l \cdot (g \cdot x) = (l \cdot g) \cdot (l \cdot x) $, $l \in \Gamma$, $ g \in LG$, $ x \in N$.
 \item[(2)] $\sigma$ is $\Gamma$-invariant.
 \item[(3)] $\Phi \colon N \to L\mathfrak{g}^*$ is $\Gamma$-equivariant.
 \end{itemize}

 We say that two $\Gamma$-compatible Hamiltonian $LG$-spaces are isomorphic
 if there exists a~$\Gamma$\nobreakdash-equi\-variant isomorphism of Hamiltonian $LG$-spaces.
\end{Definition}

\subsection[Fixed points of a Gamma-compatible space]{Fixed points of a~$\boldsymbol{\Gamma}$-compatible space}

Let $(N,\sigma,\Phi)$ be a~$\Gamma$-compatible Hamiltonian $LG$-space.
Then, we obtain a~Hamiltonian $L^{(\kappa)}G$-space from a~$\Gamma$-compatible Hamiltonian $LG$-space by taking fixed point objects.
Now $\Gamma$-fixed objects are defined by
\begin{itemize}\itemsep=0pt
 \item $N^{\Gamma} := \{x \in N \mid l \cdot x = x,\ l \in \Gamma\}$,
 \item $\sigma^{\Gamma} := \sigma|_{N^{\Gamma}}$,
 \item $\Phi ^{\Gamma} \colon N^{\Gamma} \to L\mathfrak{g}^{*}$, $x \mapsto \Phi(x)$.
\end{itemize}
In this paper, we denote $\Gamma$-fixed objects as above. For example, the $\Gamma$-fixed subgroup of~$LG$ is denoted by $LG^{\Gamma}$.
Each connected component of $N^{\Gamma}$ is a~Banach submanifold of $N$.
By considering left transformation of $LG^{\Gamma}$,
the subgroup $LG^{\Gamma}$ is a~Banach Lie subgroup of~$LG$ (see Appendix~\ref{Appendix C}).

We thus obtain the following proposition.

\begin{Proposition}
 For any connected component $L$ of $N^{\Gamma}$, the tuple
 $\bigl(L, \sigma^{\Gamma}|_L, \Phi^{\Gamma}|_L\bigr)$ is a~Hamiltonian $LG^{\Gamma}$-space, i.e., the following properties hold:
\begin{itemize}\itemsep=0pt
 \item $LG^{\Gamma}$ acts on $L$.
 \item $\sigma^{\Gamma}|_L$ is a~$LG^{\Gamma}$-invariant symplectic form on~$L$.
 \item $\Phi^{\Gamma}|_L$ is $LG^{\Gamma}$-equivariant and $\sigma^{\Gamma}|_L$, $\Phi^{\Gamma}|_L$ satisfy the moment map condition
 \[\iota\bigl(\xi^{\#}\bigr)\sigma^{\Gamma} = - {\rm d} \oint \bigl(\Phi^{\Gamma},\xi\bigr), \qquad \xi \in L\mathfrak{g}^{\Gamma}.\]
 \end{itemize}
\end{Proposition}

\begin{proof}
 Since $G$ is simply connected, the space $LG^{\Gamma}$ is connected.
 Therefore, $LG^{\Gamma}$ acts on a~connected component $L$.
 The moment map condition for $\Phi^{\Gamma}|_L$ is a straightforward calculation and
 the weak non-degeneracy of $\sigma^{\Gamma}$ is proved by averaging over $\Gamma$.
 Let $\iota \colon N^{\Gamma} \to N$ be the inclusion map.
 For any $\xi \in L\mathfrak{g}^{\Gamma}$,
 \begin{align*}
 \iota\bigl(\xi^{\#}\bigr)\sigma^{\Gamma}|_L &= \bigl(\iota\bigl(\xi^{\#}\bigr)\sigma\bigr)|_{L} = \left. -{\rm d} \oint (\Phi,\xi) \right|_{L} = - \oint \bigl(\Phi^{\Gamma} |_{L},\xi \bigr).
 \end{align*}
 Next, we will show the non-degeneracy of $\sigma^{\Gamma}$.
 Fix an element $x \in N^{\Gamma}$ and $X \in \ker\sigma^{\Gamma}$.
 Then for any $Y \in T_xN$,
 \begin{align*}
 \iota(Y)\iota(X) \sigma &= \sigma(X,Y)= \frac{1}{|\Gamma|}\sum_{l \in \Gamma} l^*\sigma(X,Y) = \sigma\left( X,\ \frac{1}{|\Gamma|} \sum_{l \in \Gamma}l_*Y \right)
 = 0.
 \end{align*}
 Thus we have proved the non-degeneracy condition.
\end{proof}

\begin{Remark}Suppose that $\Gamma$ is a~cyclic group $\bZ/m\bZ$ and that the action of $\Gamma$ on $LG$ is the one defined in Section~\ref{Twisted loop groups}.
 Then we can identify $LG^{\Gamma}$ with the twisted loop group $L^{(\kappa)}G$,
\begin{align*}
 LG^{\Gamma} &= \{g \in LG\mid l\cdot g = g,\, l \in \Gamma\} = \bigl\{g \in LG\mid \kappa g(t) =g(t+1) \text{ for all }t \in S^1_m\bigr\}
 = L^{(\kappa)}G.
\end{align*}
For any $\xi \in L^{(\kappa)}\mathfrak{g} = L\mathfrak{g}^{\Gamma}$,
\begin{align*}
 \iota\bigl(\xi^{\#}\bigr)\sigma^{\Gamma}
 = -{\rm d} \oint \bigl(\Phi^{\Gamma}, \xi\bigr) = -{\rm d} \sum_{i=1}^{m} \int_{i}^{i+1} \bigl(\Phi^{\Gamma}, \xi\bigr)
 = -m{\rm d} \int_{0}^{1} (\Phi,\xi).
\end{align*}
Therefore, $\big(L, \frac{1}{m}\sigma^{\Gamma}, \Phi^{(\kappa)}\big)$ becomes a~Hamiltonian $L^{(\kappa)}G$-space.
\end{Remark}

\subsection{Equivalence theorem}
It is known that there exists a~natural correspondence between (twisted) Hamiltonian $LG$-spaces and (twisted) quasi-Hamiltonian $G$-spaces
which is called the equivalence theorem~\cite{AMM,LYM_twisted_LG,Meinrenken2017}.
Taking base points $\beta$ with $\beta = 0 \cdot \Gamma \subset S^{1}_m$, the equivalence theorem between $\Gamma$-compatible Hamiltonian $LG$-spaces and $\Gamma$-compatible quasi-Hamiltonian $G^{\beta}$-spaces can be shown.
Let $\Omega^{\beta}G$ be a~based loop group for $\beta $, i.e., $\Omega ^{\beta}G = \{g \in LG\mid g|_{\beta} = 1_G\}$.

\begin{Theorem}\label{equiv thm}
Taking the holonomy manifold for $\Omega^{\beta}G$ yields a~bijective correspondence between
\begin{itemize}\itemsep=0pt
 \item the isomorphism class of $\Gamma$-compatible Hamiltonian $LG$-spaces $(N,\sigma,\Phi)$ with a~proper moment map,
 \item the isomorphism class of $\Gamma$-compatible $q$-Hamiltonian $G^{\beta}$-spaces $(M,\omega,\mu)$.
\end{itemize}
\end{Theorem}

\section{Twisted connections and finite group actions}\label{5Tw_conn_and_fin_grp_act}

\subsection{Preliminaries}
Let $\Sigma$ be a~compact connected 2-manifold with boundary components $V_0, V_1, \dots, V_{r}$ and $P := \Sigma \times G$ a~trivial bundle on $\Sigma$.
Fix $\lambda > 1$ and let $\Omega^{j}(\Sigma,\mathfrak{g})$ be the space of $\mathfrak{g}$-valued j-forms of~Sobolev class~${\lambda-j}$.
Then the gauge $\mathcal{G}(\Sigma) := \Map(\Sigma,G)$ is a~Banach Lie group modeled on $\Lie(\mathcal{G}(\Sigma)) = \Omega^{0}(\Sigma,\mathfrak{g})$.
Let $\mathcal{A}(\Sigma) := \Omega^{1}(\Sigma,\mathfrak{g})$ be the space of connections on the trivial bundle on $\Sigma$.
$\mathcal{A}(\Sigma)$ has a~smooth $\mathcal{G}(\Sigma)$-action given by
\smash{$ g\cdot A = g A g^{-1} + g^{*}\overline{\theta}$}
The space $\mathcal{A}(\Sigma)$ has a~canonical symplectic form called Atiyah--Bott $2$-form
$\sigma = \frac{1}{2} \int_{\Sigma} (\delta A\wedge \delta A )$
for which $\mathcal{G}(\Sigma)$-action is Hamiltonian with moment map
\[\langle \Psi(A), \xi\rangle = \int_{\Sigma} (\curv A,\xi) - \int _{\partial \Sigma} (A,\xi), \qquad \xi \in \Omega^{0}(\Sigma,\mathfrak{g}),
\]
where $\curv(A) $ is the curvature of $A$.
The space of flat connections is denoted by $\mathcal{A}_{\mathrm{flat}}(\Sigma)$.

\subsection{Finite group actions}
Let $\Gamma$ be a~finite group.
Fix a~free right $\Gamma$-action on $\Sigma$ which preserves the orientation of $\Sigma$, and a~left $\Gamma$-action on G which preserves the inner product on $\mathfrak{g}$.
Then a~right $\Gamma$-action on the trivial bundle is defined by
\[(x,a) \cdot \varphi = \bigl(x\cdot \varphi, \varphi ^{-1} \cdot a\bigr) ,\qquad (x,a) \in P,\qquad \varphi \in \Gamma.\]
This action satisfies $\Gamma$-compatibility
\[ (p\cdot g)\cdot \varphi = (p \cdot \varphi)\cdot (g \cdot \varphi), \qquad p \in P,\qquad g \in G,\qquad \varphi \in \Gamma,
\]
where $g\cdot \varphi = \varphi^{-1} \cdot g$.
We can also define left $\Gamma$-actions on $\Omega^{j}(P,\mathfrak{g})$ by
\[
\varphi \cdot \alpha = \varphi \circ \alpha \circ d\varphi , \qquad \alpha \in \Omega^{j}(P,\mathfrak{g}),\qquad \varphi \in \Gamma,
\]
and left $\Gamma$-actions on $\Omega^{j}(\Sigma,\mathfrak{g})$ can be defined in the same way.

\begin{Lemma}\label{lemm of compati for conn}
 The canonical $2$-form $\sigma$ is $\Gamma$-invariant and the curvature $\curv \colon \mathcal{A}(\Sigma) \to \Omega^{2}(\Sigma,\mathfrak{g})$ is $\Gamma$-equivariant.
 Moreover, the following $\Gamma$-compatibility holds:
 \[
 \varphi \cdot (g \cdot A) = (\varphi \cdot g) \cdot (\varphi \cdot A), \qquad \varphi \in \Gamma,\qquad g \in \mathcal{G}(\Sigma),\qquad A \in \mathcal{A}(\Sigma).
 \]
\end{Lemma}

\begin{proof}For any $\varphi \in \Gamma$ and $X, Y \in T_A\mathcal{A}(\Sigma) = \Omega^{1}(\Sigma,\mathfrak{g})$, the following holds:
 \[\varphi^*\sigma(X,Y) = \int_{\Sigma} (\varphi X \circ {\rm d} \varphi \wedge \varphi Y \circ {\rm d} \varphi) = \int_{\Sigma}(X\wedge Y) = \sigma(X,Y).\]

Since the exterior derivative ${\rm d} \colon \Omega^{j}(\Sigma,\mathfrak{g}) \to \Omega^{j+1}(\Sigma,\mathfrak{g})$ is $\Gamma$-equivariant, the $\Gamma$-equivariance of the curvature holds
 \begin{align*}
 \curv(\varphi \cdot A) &= {\rm d}(\varphi \cdot A) + \frac{1}{2}[ \varphi \cdot A \wedge \varphi \cdot A] = \varphi \cdot {\rm d} A + \frac{1}{2}\varphi \cdot [A \wedge A]
 = \varphi \cdot \curv A,
 \end{align*}
 where $\varphi \in \Gamma$ and $A \in \mathcal{A}(\Sigma)$.
 Moreover, for any $g \in \mathcal{G}(\Sigma)$, $\varphi \in \Gamma$ and $A \in \mathcal{A}(\Sigma)$, the following holds
$\varphi \circ \Ad(g) = \Ad(\varphi \cdot g) \circ \varphi$,
$\varphi \cdot g^*\overline{\theta} = (\varphi \cdot g)^* \overline{\theta}$.
 Therefore, we obtain the $\Gamma$-compatibility of~$\mathcal{G}(\Sigma)$ and $\mathcal{A}(\Sigma)$.
\end{proof}

$\Gamma$-actions on $\mathcal{G}(\partial \Sigma)$ and $\Omega^{j}(\partial \Sigma, \mathfrak{g})$ are also defined in the same way.

\subsection[The space M(Sigma) and Gamma-compatibilities]{The space $\boldsymbol{\mathcal{M}(\Sigma)}$ and $\boldsymbol{\Gamma}$-compatibilities}

Since the restriction map $\mathcal{G}(\Sigma) \to \mathcal{G}(\partial \Sigma)$ is surjective,
the following sequence is exact:
\[ 1 \to \mathcal{G}_{\partial \Sigma}(\Sigma) \to \mathcal{G}(\Sigma) \to \mathcal{G}(\partial \Sigma) \to 1,
\]
where $\mathcal{G}_{\partial \Sigma}(\Sigma) := \{g \in \mathcal{G}(\Sigma)\mid g|_{\partial \Sigma} =1_G\}$.
Then $\mathcal{M}(\Sigma) := \mathcal{A}_{\mathrm{flat}}(\Sigma) / \mathcal{G}_{\partial \Sigma}(\Sigma)$ is a~Banach manifold and a~Hamiltonian $(\mathcal{G}(\partial \Sigma) \simeq LG^{r+1})$-space with proper moment map~\cite{AMM,Donaldson1992}.
The symplectic form~$\sigma$ on $\mathcal{M}(\Sigma)$ is induced from Atiyah--Bott $2$-form on $\mathcal{A}(\Sigma)$ and the moment map is a~restriction
\[
\Phi \colon \mathcal{M}(\Sigma) \to \Omega^1(\partial \Sigma, \mathfrak{g}),\qquad A \mapsto A|_{\partial \Sigma}.
\]
The $\Gamma$-compatibility of $\mathcal{M}(\Sigma)$ can be shown by using Lemma~\ref{lemm of compati for conn}. A~$\Gamma$-action on $\mathcal{M}(\Sigma)$ is defined~by
$\varphi \cdot [A] := [\varphi \cdot A]$, $ \varphi \in \Gamma$, $ [A] \in \mathcal{M}(\Sigma)$.
The Hamiltonian $\mathcal{G}(\partial \Sigma)$-space $(\mathcal{M}(\Sigma), \sigma, \Phi)$ together with the $\Gamma$-action is a~$\Gamma$-compatible Hamiltonian $\mathcal{G}(\partial \Sigma)$-space.

\begin{Proposition}
 $(\mathcal{M}(\Sigma),\sigma,\Phi)$ is a~$\Gamma$-compatible Hamiltonian $\mathcal{G}(\partial \Sigma)$-space.
\end{Proposition}
The above proposition follows immediately from Lemma~\ref{lemm of compati for conn}.

\subsection[Fixed points of M(Sigma)]{Fixed points of $\boldsymbol{\mathcal{M}(\Sigma)}$}

By taking a~fixed point objects $\bigl(\mathcal{M}(\Sigma)^{\Gamma},\omega^{\Gamma},\Phi^{\Gamma}\bigr)$, we get Hamiltonian spaces for twisted loop groups actions.

We define a~$\Gamma$-action on $I := \{ 0,1,\dots,r\}$ by using a~permutation of $\partial \Sigma = \coprod _{i=0}^{r} V_i$,
$V_{i \cdot \varphi } = V_{i} \cdot \varphi$.
Let $\Gamma(i) $ be a~stabilizer of $i \in I$, i.e., $\Gamma(i) = \{\varphi \in \Gamma\mid i \cdot \varphi = i\}$.

\begin{Lemma}
 For any $i \in I$, $\Gamma(i)$ is a~cyclic group.
\end{Lemma}

\begin{proof}
 Fix a~point $b \in V_i$ and consider a~$\Gamma(i)$-orbit $b \cdot \Gamma(i)$, we restrict the $\Gamma(i)$-action on $\partial \Sigma$ to~$b \cdot \Gamma(i)$.
 Since this $\Gamma$-action on $\partial \Sigma$ is free, the restricted $\Gamma(i) $-action on $b \cdot \Gamma(i)$ is simply transitive.
 So by fixing a~cyclic order of $b \cdot \Gamma(i) =\{b_0, b_1, \dots, b_{m_i-1}\}$, we get a~unique element~$b_1 = b_0 \cdot \kappa_i$.
 Then, the element $\kappa_i$ is a~generator of $\Gamma(i) = \langle \kappa_i \rangle$.
\end{proof}

We fix a~generator $\kappa_i$ of $\Gamma(i)$ for all $i \in I$.
We set $m_i := | \Gamma(i)|$, and identify $V_i$ with $S^{1}_{m_i} = \bR / m_i \bZ$ while keeping $\Gamma(i)$-action.
So, the following isomorphism is obtained:
\smash{$
 \mathcal{G}(V_i)^{\Gamma(i)} \simeq L^{(\kappa_i)}G$}.
So, by taking a~complete system of representatives $I^{\prime}$ of $I/\Gamma$, and fixing $\kappa_i$ for any $i \in I$,
and taking~${\tau_{ji} \in \Gamma }$ such that \smash{$\kappa_j = \tau_{ji} \kappa_i \tau_{ji}^{-1}$} for $i, j \in i\cdot \Gamma$,
we identify the fixed points subgroup~$\mathcal{G}(\partial \Sigma)^{\Gamma}$ with a~product of twisted loop groups,
\begin{align*}
 \mathcal{G}(\partial \Sigma)^{\Gamma} &\simeq \left\{(g_i) \in \underset{i \in I}{\prod} L^{(\kappa_i)}G\,\bigg|\, g_i = \tau_{ij} g_j\right\} \simeq \underset{i \in I^{\prime}}{\prod} L^{(\kappa_i)}G.
\end{align*}
In the same way, we obtain the following isomorphisms:
\begin{align*}
 \Omega^{1}(\partial \Sigma, \mathfrak{g})^{\Gamma} &\simeq \underset{i \in I^{\prime}}{\prod} L^{(\kappa_i)}\mathfrak{g}^{*};\qquad A \mapsto (A|_{V_i}), \qquad
 \Omega^{0}(\partial \Sigma, \mathfrak{g})^{\Gamma} \simeq \underset{i \in I^{\prime}}{\prod} L^{(\kappa_i)}\mathfrak{g};\qquad \xi \mapsto (\xi|_{V_i}).
\end{align*}
We define $\Phi^{\prime}$ and $\sigma^{\prime}$ by
\smash{$
 \Phi^{\prime} \colon \mathcal{M}(\Sigma)^{\Gamma} \to {\prod}_{i \in I^{\prime}} L^{(\kappa_i)}\mathfrak{g}^{*}$}; $ A \mapsto (\Phi(A)|_{V_i})$, \smash{$
 \sigma^{\prime} := \frac{1}{|\Gamma |} \sigma$}.
By considering the \smash{$\prod_{i \in I} L^{(\kappa_i)}G$}-action on $\mathcal{M}(X)$, the following proposition holds.

\begin{Proposition}
 Each connected component $\mathcal{N}$ of $\mathcal{M}(\Sigma)^{\Gamma}$ together with the $2$-form $\sigma^{\prime}|_{\mathcal{N}}$ and the moment map $\Phi^{\prime}|_{\mathcal{N}}$ is a~Hamiltonian \smash{$\Pi_{i \in I^{\prime}} L^{(\kappa_i)}G$}-space.
\end{Proposition}

\begin{proof}
 It is trivial that $\sigma^{\prime}$ is symplectic and $\Phi^{\prime}$ is equivariant for twisted loop group actions.
 So, we prove the moment map condition:
 for any $[A] \in \mathcal{M}(\Sigma)$ and \smash{$(\xi_i) = \xi \in \prod_{i \in I^{\prime}} L^{(\kappa_i)}\mathfrak{g}$},
\begin{align*}
 \iota\bigl(\xi^{\#}\bigr)\sigma^{\prime}_{[A]} &= \frac{1}{| \Gamma|} \iota\bigl(\xi^{\#}\bigr)\sigma_{[A]}|_{\Gamma}
 = -\frac{1}{| \Gamma|}{\rm d} \int_{\partial \Sigma} (\Phi(A), \xi)
 = -\frac{1}{| \Gamma|} \sum_{i \in I} {\rm d} \oint_{V_i}(A|_{V_i}, \xi|_{V_i}) \\
 &= -\frac{1}{| \Gamma|} \sum_{i \in I^{\prime}} \left| i \cdot \Gamma\right| {\rm d} \oint_{V_i}(A|_{V_i}, \xi_i)
 = -\frac{1}{| \Gamma|} \sum_{i \in I^{\prime}} \left| i \cdot \Gamma\right| \left| \Gamma(i)\right| {\rm d} \int_{0}^{1} (A|_{V_i}, \xi_i) \\
 &
 = - \sum_{i \in I^{\prime}} {\rm d}\int_{0}^{1} \bigl(\Phi^{\prime}(A)_i, \xi_i\bigr).\tag*{\qed}
\end{align*} \renewcommand{\qed}{}
\end{proof}

\subsection[The space M(Sigma) and Gamma-compatibilities]{The space $\boldsymbol{M(\Sigma)}$ and $\boldsymbol{\Gamma}$-compatibilities}


Let $\beta \subset \partial \Sigma$ be a~finite subset (base points) which is $\Gamma$-invariant and we fix a~cyclic ordering of~$\beta_i = \bigl(b^{i}_{\lambda}\bigr) := \beta \cap V_i$.
We denote
$\mathcal{G}_{\beta}(\Sigma) := \{g \in \mathcal{G}(\Sigma)\mid g|_{\beta} = 1\}$, $
 \mathcal{G}_{\beta}(\partial \Sigma) := \{g \in \mathcal{G}(\partial \Sigma)\mid g|_{\beta}=1\}$.
Set $M(\Sigma) := \mathcal{A}_{\mathrm{flat}}(\Sigma)/\mathcal{G}_{\beta}(\Sigma) \simeq \mathcal{M}(\Sigma)/\mathcal{G}_{\beta}(\partial \Sigma)$.
An action of $\mathcal{G}(\partial \Sigma)/\mathcal{G}_{\beta}(\partial \Sigma) \simeq G^{\beta}$ on $M(\Sigma)$ is induced from the residual gauge action on $\mathcal{M}(\Sigma)$.
The collection of holonomy map \smash{$\Hol^{b^i_{\lambda}}_{b^i_{\lambda+1}}(A)$} defines a~$G^{\beta}$-equivariant map
\[
\mu \colon\ M(\Sigma) \to \bH,\qquad A \mapsto \Hol(A) := \bigl(\Hol^{b^i_{\lambda}}_{b^i_{\lambda+1}}(A)\bigr)_{b^i_{\lambda} \in \beta},
\]
where $\bH = G^{\beta}$ is a~$G^{\beta}$-bitorsor under the following left and right actions,
$(g \cdot h)\bigl(b^{i}_{\lambda}\bigr) = g\bigl(b^i_{\lambda}\bigr)h\bigl(b^i_{\lambda}\bigr)$, $
 (h \cdot g)\bigl(b^i_{\lambda}\bigr) = h\bigl(b^i_{\lambda}\bigr)g\bigl(b^{i}_{\lambda+1}\bigr)$, $ g \in G^{\beta}$, $ h \in \bH$.
An action of $\Gamma$ on $M(\Sigma)$ is induced from the $\Gamma$-action on $\mathcal{M}(\Sigma)$.
The following theorem, which is analogous to~\cite[Theorem 9.1]{AMM}, holds.


\begin{Theorem}\label{q-Ham str on M(X)}
 There is a~natural $G^{\beta}$-invariant $2$-form $\omega$ on $M(\Sigma)$ for which $(M(\Sigma),\omega,\mu)$ is a~$\Gamma$-compatible $q$-Hamiltonian $G^{\beta}$-space.
\end{Theorem}

\begin{proof}
 The proof is analogous to that of~\cite[Theorem 9.1]{AMM}.
 In this proof, let $\pi \colon \mathcal{M}(\Sigma) \to M(\Sigma) $ be the quotient map.
 First, we will show the equivariance of $\mu \colon M(\Sigma) \to \bH$.
 For any $g \in G^{\beta}$ and $\pi(A) \in M(\Sigma)$,
$\mu(g \cdot \pi(A)) = \mu (\pi(\tilde{g}\cdot A)) = \Hol(\tilde{g}\cdot A)= \Ad_{g} \mu(\pi(A))$.

Let $R_{j} \colon \Omega^1(\partial \Sigma, \mathfrak{g}) \to \Omega^1(V_j,\mathfrak{g})$ be a~restriction and let $\varpi_{j}$ be the canonical $2$-form on $\Omega^{1}(V_j, \mathfrak{g}) \simeq L\mathfrak{g}^*$.
 Denote $\varpi := \sum_{j=0}^{r} R_{j}^*\varpi_{j}$. We define a~$2$-form on $\mathcal{M}(\Sigma)$, $\tilde{\omega} := \sigma - \Phi^*\varpi$.
 Then the $2$-form $\tilde{\omega}$ is basic for $\mathcal{G}_{\beta}(\partial \Sigma)$-action for $\mathcal{M}(\Sigma)$.
 Since $\sigma$ and $\varpi$ are $\mathcal{G}(\partial \Sigma)$-invariant, $\tilde{\omega}$~becomes a~$\mathcal{G}(\partial \Sigma)$-invariant $2$-form.
 Furthermore, for any $\xi \in \Lie(\mathcal{G}_{\beta}(\partial \Sigma))$,
 \begin{align*}
 \iota\bigl(\xi^{\#}\bigr)\tilde{\omega}_A &= \iota\bigl(\xi^{\#}\bigr)\sigma_A - \iota\bigl(\xi^{\#}\bigr) \Phi^* \varpi_A
 = \iota\bigl(\xi^{\#}\bigr) \sigma_A - \sum_{j=0}^{r} {R_j}^* \bigl(\iota\bigl({\xi|_{V_j}}^{\#}\bigr)\varpi_j\bigr)\\
 &= -{\rm d} \int_{\partial \Sigma }(A|_{\partial \Sigma}, \xi)- \sum_{j=0}^{r} {R_{j}}^*\left(-{\rm d} \oint (A|_{V_j},\xi|_{V_j}) -\frac{1}{2}(\Hol_{j})^* \bigl(\theta_{j}+ \overline{\theta}_{j}, \xi|_{\beta_{j}}\bigr) \right)= 0,
 \end{align*}
 where $\Hol_j \colon \Omega^1(V_{j},\mathfrak{g}) \to \bH^{j} := G^{\beta_j}$, \smash{$A_j \mapsto \bigl(\Hol^{b^{j}_\lambda}_{b^{j}_\lambda}(A_j)\bigr)$}, and $\theta_j $ (resp.\ $\overline{\theta}$) is the left (resp.\ right) Maurer--Cartan form on~$\bH^j$.
 So, $\tilde{\omega}$ induces a~unique $G^{\beta}$-invariant $2$-form $\omega$ on $M(\Sigma)$ which satisfies $\pi^* \omega = \tilde{\omega}$.

 Now we prove the tuple $(M(\Sigma),\omega,\mu)$ is a~twisted quasi-Hamiltonian $G^{\beta}$-space.
 We denote the canonical $3$-form on $\bH$ by $\chi_{\bH}$, i.e., $\chi_{\bH} = \frac{1}{12}(\theta_{\bH},[\theta_{\bH},\theta_{\bH}])$.
 Then, by $\Hol \circ \Phi = \mu \circ \pi$,
$
 \pi^*{\rm d}\omega = {\rm d}\tilde{\omega}= {\rm d}( \sigma - \Phi^* \varpi)
 = \Phi^* \Hol^* \chi_{\bH}
 = \pi^* \mu^* \chi_{\bH}$.
 Therefore, we have proved ${\rm d} \omega = \mu^* \chi_{\bH}$.

 Next, we want to show the moment map condition.
 For any $\xi \in \mathfrak{g}^{\beta}$, there exists $\tilde{\xi} \in \Omega^0(\partial \Sigma,\mathfrak{g})$ such that $\tilde{\xi}|_{\beta} = \xi$.
 Then
 \begin{align*}
 \pi^*\bigl(\iota\bigl(\xi^{\#}\bigr)\omega\bigr) &= \iota\bigl(\tilde{\xi}\bigr) \pi^*\omega
 = \iota\bigl(\tilde{\xi}^{\#}\bigr) \bigl(\sigma-\Phi^*\varpi\bigr)
 = \frac{1}{2} \sum_{j=0}^{r} {R_j}^*{\Hol_j}^*\bigl(\theta_j+ \overline{\theta}_j,\tilde{\xi}|_{\beta_{j}}\bigr)\\
 &= \pi^* \left( \frac{1}{2} \mu^* \bigl(\theta_{\bH} + \overline{\theta}_{\bH},\xi\bigr)\right).
 \end{align*}
 So, the moment map condition holds.

 Finally, we will prove the degeneracy of $\omega$, i.e., we will show $\ker \omega_m \cap \ker d_m \mu = \{0\}$ for any~$m \in M(\Sigma)$.
 Fix $v \in \ker \omega_m \cap \ker {\rm d}_m \mu$ and $A \in \mathcal{M}(\Sigma)$ such that $\pi(A)=m$.
 Since $\pi \colon \mathcal{M}(\Sigma) \to M(\Sigma)$ is a~submersion, there exists a~tangent vector $X \in T_A \mathcal{M}(\Sigma)$ such that ${\rm d}_A\pi(X)=v$.

 For any $Y \in T_A\mathcal{M}(\Sigma)$,
$\tilde{\omega}_A(X,Y) = \pi^*\omega _A(X,Y) = 0$.
 Therefore, $X \in \ker \tilde{\omega}_A$. By the definition of $\tilde{\omega}$, the following holds:
$\iota(X) \sigma= \iota(X) \Phi^*\varpi$.
Thus, $X \in (\ker {\rm d}_A \Phi )^{\sigma}$.
 Since the space $( \ker {\rm d}_A \Phi )^{\sigma}$ equals to the tangent space of the orbit $\mathcal{G}(\partial \Sigma) \cdot A$, i.e., there exists $\xi \in \Omega^0(\partial \Sigma, \mathfrak{g})$ that the generating vector $\xi^{\#}$ equals to~$X$ \cite[Theorem~9.1]{AMM}.

 Since \smash{${\xi|_{\beta}}^{\#} \in \ker {\rm d}_m \mu$}, we obtain
 $\xi|_{\beta} \in \ker(\Ad_{\mu(m)}-1)$.
 Now we define $\eta_j \in \Omega^0(V_j,\mathfrak{g})$ by~\smash{$
 \eta_j(s) := \Ad\bigl(\Hol^{b^{j}_0}_{s}(\Phi(A))^{-1}\bigr)\xi\bigl(b^{j}_0\bigr)$}.
 Because \smash{$\Ad_{\mu(m)}\xi|_{\beta} = \xi|_{\beta}$},
 \[
 \eta_j\bigl(b^{j}_{\lambda_j}\bigr) = \Ad\bigl(\Hol^{b^{j}_0}_{b^{j}_{\lambda_j}}(\Phi(A))\bigr)\xi\bigl(b^{j}_0\bigr) = \xi\bigl(b^{j}_{\lambda_j}\bigr).
 \]
 Therefore, $\eta_j$ is well defined.
 By collecting $\eta_j$ for all $j =0,\dots,r$, we define $\eta \in \Omega^0(\partial \Sigma, \mathfrak{g})$.
 Because $(\eta-\xi)|_{\beta} = 0$, $\eta^{\#} = (\eta - \xi)^{\#} + \xi^{\#}$ is in $\ker\bigl(\sigma- \Phi^* \varpi\bigr)$.

We now prove $\iota\bigl(\eta^{\#}\bigr)\Phi^*\varpi=0$.
 For any $\alpha \in T_A \mathcal{M}(\Sigma)$, by Proposition~\ref{formula on loop group} and Lemma~\ref{lemma of calc of Hol},
 \begin{align*}
 \iota(\alpha)\iota\bigl(\eta^{\#}\bigr) \Phi^*\varpi &= \sum_{j=0}^{r} \iota(\alpha|_{V_j}) \iota\bigl({\eta|_{V_j}}^{\#}\bigr) \varpi_j\\
 &=-\sum_{j=0}^{r} \iota(\alpha|_{V_j}) {\rm d} \oint (A|_{V_j}, \eta|_{V_j}) - \frac{1}{2} \iota(\alpha)\Hol^*\bigl(\theta_{\bH}+ \overline{\theta}_{\bH}, \eta|_{\beta}\bigr) \\
 &=-\sum_{j=0}^{r} \oint (\alpha|_{V_j}, \eta|_{V_j}) - \iota(\alpha) \iota\bigl(\xi^{\#}\bigr) \omega\\
 &= -\sum_{j=0}^{r} \sum_{\lambda_j = 1}^{m_j} \int_{b^j_{\lambda_j}}^{b^j_{\lambda_{j+1}}} \bigl(\alpha|_{V_j}, \Ad\bigl(\Hol^{b^j_{\lambda_{j}}}_{s}(A|_{V_j})^{-1}\bigr)\xi\bigl(b^j_{\lambda_j}\bigr)\bigr){\rm d} s\\
 &=\sum_{j=0}^{r} \sum_{\lambda_j =1}^{m_j} \int _{b^j_{\lambda_j}}^{b^j_{\lambda_{j+1}}} \bigl(\Ad\bigl(\Hol^{b^j_{\lambda_j}}_{s}(A)\bigr)\alpha|_{V_j}, \xi\bigl(b^j_{\lambda_j}\bigr)\bigr) \\
 &= \sum_{j=0}^{r} \sum_{\lambda_j =1}^{m_j} \big( \iota(\alpha|_{V_j}){\big(\Hol^{b^j_{\lambda_j}}_{b^j_{\lambda_{j+1}}}\big)}^*\overline{\theta}, \xi\bigl(b^j_{\lambda_j}\bigr)\big) \\
 &= \frac{1}{2} \sum_{j=0}^{r} \sum_{\lambda_j=0}^{m_j} \big( \iota(\alpha|_{V_j}){\big(\Hol^{b^j_{\lambda_j}}_{b^j_{\lambda_{j+1}}}\big)}^*\overline{\theta}, \xi\bigl(b^j_{\lambda_j}\bigr) \big)\\
 & \qquad + \frac{1}{2} \sum_{j=0}^{r} \sum_{\lambda_j=0}^{m_j} \big( \iota(\alpha|_{V_j}){\big(\Hol^{b^j_{\lambda_j}}_{b^j_{\lambda_{j+1}}}\big)}^*\overline{\theta}, \Ad\bigl(\Hol^{b^j_{\lambda_j}}_{b^j_{\lambda_{j+1}}}\bigr)\xi\bigl(b^j_{\lambda_{j+1}}\bigr) \big) \\
 &= \iota(\alpha) \frac{1}{2} \Hol^*\bigl(\theta_{\bH}+ \overline{\theta}_{\bH}, \xi|_{\beta}\bigr) =0.
 \end{align*}
 Considering \smash{$\iota\bigl(\eta^{\#}\bigr) \Phi^*\varpi=0$} and \smash{$\eta^{\#} \in \ker\bigl(\sigma- \Phi^*\varpi\bigr)$}, we obtain $\smash{\iota\bigl(\eta^{\#}\bigr)\sigma= \iota\bigl(\eta^{\#}\bigr)}\bigl(\bigl(\sigma - \Phi^*\varpi\bigr)+\smash{\Phi^*\varpi}\bigr) =0$.
 Since $\sigma$ is weak non-degenerate, \smash{$\eta^{\#}$} is zero.
 Therefore, we conclude $\smash{{\xi|_{\beta}}^{\#} }= \smash{{(\xi-\eta)|_{\beta}}^{\#} + {\eta|_{\beta}}^{\#}}=0$.
\end{proof}

\begin{Remark}\quad
\begin{itemize}\itemsep=0pt
 \item[(1)] If the group $\Gamma$ is trivial, we can choose any finite subset of $\partial \Sigma$ as the set of base points.
 Then the space $M(\Sigma)$ is a~straightforward generalization of the space in~\cite[Theorem~9.1]{AMM}.
 \item[(2)] Although Theorem~\ref{q-Ham str on M(X)} follows directly from Theorem~\ref{equiv thm}, we include a~proof because the almost same argument applies to the case where $G$ is not simply connected.
 \end{itemize}
\end{Remark}

\subsection[Fixed points of M(X)]{Fixed points of $\boldsymbol{M(X)}$}\label{section_twisted_reps}
Taking fixed point objects of $(M(\Sigma),\omega,\mu)$, we obtain a~twisted $q$-Hamiltonian $\bigl(G^{\beta}\bigr)^{\Gamma}$-space $\bigl(M(\Sigma)^{\Gamma},\omega^{\Gamma},\mu^{\Gamma}\bigr)$.
We denote $ \underline{\Sigma} := \Sigma /\Gamma$ and $\beta_{\underline{\Sigma}} := \pi(\beta)$, where $\pi \colon \Sigma \to \underline{\Sigma}$ is the covering map.

Taking a~holonomy yields the following correspondence: $M(\Sigma) \simeq \Hom(\Pi_1(\Sigma,\beta),G)$.
This correspondence is $\Gamma$-equivariant for an action on $\Hom(\Pi_1(\Sigma,\beta),G)$
$(\varphi \cdot \rho)(\gamma) = \varphi \rho(\gamma\cdot\varphi)$, $
 \varphi \in \Gamma$, ${\rho \in \Hom(\Pi_1(\Sigma,\beta),G)}$, $ \gamma \in \Pi_1(\Sigma,\beta)$.
Therefore, we obtain a~bijection $M(\Sigma)^{\Gamma} \simeq \Hom(\Pi_1(\Sigma,\beta),G)^{\Gamma}$.

Now we want to consider the space $\Hom(\Pi_1(\Sigma,\beta),G)^{\Gamma}$ as the space on $\underline{\Sigma}$.
Let $G_{\underline{\Sigma}}$ be a~group bundle $(\Sigma \times G) /\Gamma$ on $\underline{\Sigma}$ with a~fiber $G$.
By fixing a~complete system of representatives $I$ of $\beta$ for~$\Gamma$-action,
we obtain the monodromy representation of $G_{\underline{\Sigma}}$ denoted by $\mon_{\underline{\Sigma},I} \colon \Pi_1(\underline{\Sigma},\beta_{\underline{\Sigma}}) \to \Gamma$ (cf.\ Appendix~\ref{Appendix B}).
We denote the space of representations twisted by the morphism $\mon_{\underline{\Sigma},I}$ by
\[
\Hom_{\mon_{\underline{\Sigma},I}}(\Pi_1(\underline{\Sigma},\beta_{\underline{\Sigma}}),G):= \bigr\{\underline{\rho} \in \Hom(\Pi_1(\underline{\Sigma},\beta_{\underline{\Sigma}}),\Gamma \ltimes G)\mid \pr \circ \underline{\rho} = \mon_{\underline{\Sigma},I}\bigr\}.
\]
Then, the covering map $\pi \colon \Sigma \to \underline{\Sigma}$ induces a~bijection
\[
\Hom_{\mon_{\underline{\Sigma},I}}(\Pi_1(\underline{\Sigma},\beta_{\underline{\Sigma}}),G) \simeq \Hom(\Pi_1(\Sigma,\beta),G)^{\Gamma}.
\]
Therefore, we obtain a~bijection
$M(\Sigma)^{\Gamma} \simeq \Hom_{\mon_{\underline{\Sigma},I}}(\Pi_1(\underline{\Sigma},\beta_{\underline{\Sigma}}),G)$.
By choosing a~system of free generators of $\Pi_1(\underline{\Sigma},\beta_{\underline{\Sigma}})$, the space $\Hom_{\mon_{\underline{\Sigma},I}}(\Pi_1(\underline{\Sigma},\beta_{\underline{\Sigma}}),G)$ is diffeomorphic to a~product of $G$-bitorsors.
Since $G$ is connected, this product of $G$-bitorsors is connected.
Therefore, the tuple~${\bigl(M(\Sigma)^{\Gamma}, \omega^{\Gamma}, \mu^{\Gamma}\bigr)}$ is a~twisted quasi-Hamiltonian \smash{$\bigl(G^{\beta}\bigr)^{\Gamma}$}-space.
Also \smash{$\Hom_{\mon_{\underline{\Sigma},I}}(\Pi_1(\underline{\Sigma},\beta_{\underline{\Sigma}}),G)$} carries a~twisted quasi-Hamiltonian $G^{\beta_{\underline{\Sigma}}}$-structure.
The above discussion is summarized in the following theorem.

\begin{Theorem}\label{q-Ham str on rep sp}
 The representation space $\Hom(\Pi_1(\Sigma,\beta),G)$ has a~structure of a~$\Gamma$-compatible quasi-Hamiltonian $G^{\beta}$-space.
 By taking the set of $\Gamma$-fixed points of $\Hom(\Pi_1(\Sigma,\beta),G)$, we obtain a~structure of a~twisted quasi-Hamiltonian $G^{\beta_{\underline{\Sigma}}}$-space on $\Hom_{\mon_{\underline{\Sigma}, I}}(\Pi_1(\underline{\Sigma},\beta_{\underline{\Sigma}}),G)$.
\end{Theorem}


\section{Twisted group bundles and monodromy}\label{Appendix B}

\subsection{Setting}

Let $\Sigma$ be an oriented compact 2-manifold with boundary.
Let $\Gamma$ be a~finite group.
We consider a~left $\Gamma$-action on $G$ and a~free right $\Gamma$-action on $\Sigma$.
We define
 $ \underline{\Sigma} := \Sigma/\Gamma$,
 trivial group bundle $G_{\Sigma} := \Sigma \times G$ on $\Sigma$,
a~group bundle $G_{\underline{\Sigma}}:= (\Sigma \times G)/\Gamma$ on $\underline{\Sigma}$.

Let $\beta_{\underline{\Sigma}} $ be a~finite subset of $\underline{\Sigma}$, and let $\beta_{\Sigma}$ be a~pre-image of $\beta_{\underline{\Sigma}}$ for the quotient map $q \colon \Sigma \to \underline{\Sigma}$.
We denote the fundamental groupoid by $\Pi_{\Sigma} := \Pi_1(\Sigma,\beta_{\Sigma})$ and $\Pi_{\underline{\Sigma}} := \Pi_1(\underline{\Sigma},\beta_{\underline{\Sigma}})$.

\subsection[The monodromy of G\_Sigma]{The monodromy of $\boldsymbol{G_{\underline{\Sigma}}}$}
The trivial group bundle $G_{\Sigma}$ has the trivial flat Ehresmann connection $\mathcal{H}_{\Sigma}$. We denote its monodromy by $\mon_{\Sigma}$.
Since $G_{\Sigma}$ is trivial as a~group bundle, under the natural identification~${G_b \simeq G}$ for all $b \in \beta_{\Sigma}$, we obtain the trivial monodromy representation
$
\mon_{\Sigma, \beta} \colon\Pi_{\Sigma} \to \Aut(G), \gamma \mapsto 1$.

Since $\mathcal{H}_{\Sigma}$ is $\Gamma$-invariant under the diagonal action of $\Gamma$ on $G_{\Sigma}$, $\mathcal{H}_{\Sigma}$ induces a~flat Ehresmann connection $\mathcal{H}_{\underline{\Sigma}}$ on $\underline{\Sigma}$.
We denote its monodromy by $\mon_{\underline{\Sigma}}$.
For any $\gamma \in \Pi_{\underline{\Sigma}}$, $\mon_{\underline{\Sigma}}(\gamma) \in \Isom(G_{s(\gamma)},G_{t(\gamma)})$.

Fix a~complete system of representatives $I \subset \beta_{\Sigma}$ for the $\Gamma$-action.
We construct a~$\Gamma$-valued representation of $\mon_{\underline{\Sigma}}$.
For any $\gamma \in \Pi_{\underline{\Sigma}}$, $i(s(\gamma)) \in I$ denotes a~representative of~$q^{-1}(s(\gamma))$, and~${i(t(\gamma)) \in I}$ denotes a~representative of $q^{-1}(t(\gamma))$.
Let $\tilde{\gamma}$ be a~lift of $\gamma \in \Pi_{\underline{\Sigma}}$ starting at~$i(s(\gamma))$,
then there exists a~unique element $\varphi_I(\gamma) \in \Gamma$ such that
$t(\tilde{\gamma}) = i(t(\gamma))\cdot \varphi_I(\gamma)$.
Set ${\mon_{\underline{\Sigma},I}(\gamma) := \varphi_I(\gamma)}$,
we obtain the monodromy representation $\mon_{\underline{\Sigma},I} \colon \Pi_{\underline{\Sigma}} \to \Gamma$.

\subsection[The G-valued representation of Pi\_Sigma]{The $\boldsymbol{G}$-valued representation of $\boldsymbol{\Pi_{\underline{\Sigma}}}$}

We denote
\[
\Hom_{\mon_{\underline{\Sigma},I}}(\Pi_{\underline{\Sigma}},G) := \{\underline{\rho} \in \Hom(\Pi_{\underline{\Sigma}}, \Gamma \ltimes G)\mid \pr_{\Gamma} \circ \underline{\rho} = \mon_{\underline{\Sigma},I}\}.
\]
In this appendix, we construct a~bijection
\[
\Hom(\Pi_{\Sigma},G)^{\Gamma} \to \Hom_{\mon_{\underline{\Sigma},I}}(\Pi_{\underline{\Sigma}},G), \qquad \rho \mapsto \underline{\rho}_I.
\]
 A~$\Gamma$-action on $\Hom(\Pi_{\Sigma},G)$ is defined by
$
(\varphi \cdot \rho)(\tilde{\gamma}) = \varphi \rho(\tilde{\gamma}\cdot\varphi) $, $ \rho \in \Hom(\Pi_{\Sigma},G)$, $ \varphi \in \Gamma$, $ \tilde{\gamma} \in \Pi_{\Sigma}$.
For any $\rho \in \Hom(\Pi_{\Sigma},G)^{\Gamma}$ and $\gamma \in \Pi_{\underline{\Sigma}}$, we denote the lifting of $\gamma$ starting at $i(s(\gamma))$ by $\tilde{\gamma}$.
We define $\underline{\rho}_{I}(\gamma) := (\mon_{\underline{\Sigma},I}(\gamma),\rho(\tilde{\gamma}))$.

\begin{Proposition}
 \smash{$\underline{\rho}_I \colon \Pi_{\underline{\Sigma}} \to \Gamma \ltimes G$} is a~groupoid homomorphism.
\end{Proposition}

\begin{proof}
 For any $\gamma_1,\gamma_2 \in \Pi_{\underline{\Sigma}}$ satisfying $s(\gamma_1)=t(\gamma_2)$, \smash{$\widetilde{(\gamma_1 \gamma_2)} = (\tilde{\gamma_1}\cdot \mon_{Y,I}(\gamma_2))\tilde{\gamma_2}$}.
 Therefore,
 \begin{align*}
 \underline{\rho}_I(\gamma_1\gamma_2) &= (\mon_{Y,I}(\gamma_1 \gamma_2), \rho(\widetilde{\gamma_1 \gamma_2}))
 = (\mon_{Y,I}(\gamma_1) \mon_{Y,I}(\gamma_2), \rho(\tilde{\gamma_1}\cdot\mon_{Y,I}(\gamma_2))\rho(\tilde{\gamma_2})) \\
 &= (\mon_{Y,I}(\gamma_1) \mon_{Y,I}(\gamma_2), (\rho(\tilde{\gamma_1}) \cdot \mon_{Y,I}(\gamma_2)) \rho(\tilde{\gamma_2}))
 = \underline{\rho}_{I}(\gamma_1) \underline{\rho}_{I}(\gamma_2).\tag*{\qed}
 \end{align*} \renewcommand{\qed}{}
\end{proof}

It is easy to see that the above correspondence is bijective.

\section[A C\string^infty structure of fixed point sets]{A $\boldsymbol{C^{\infty}}$ structure of fixed point sets}\label{Appendix C}

In this appendix, we show that each connected component of the fixed point set $N^{\Gamma}$ for a~$C^{\infty}$-action of a~finite group $\Gamma$ on a~Banach manifold $N$ is a~Banach submanifold of $N$.

Fix a~point $x \in N^{\Gamma}$, a~$\Gamma$-action on the tangent space $T_xN$ is induced by the $\Gamma$-action on $N$.
We denote the averaging by $\ave_x \colon T_xN \to (T_xN)^{\Gamma}$, $ X \mapsto \frac{1}{|\Gamma|} \sum_{\varphi \in \Gamma} \varphi \cdot X$.
Because the $\Gamma$-action on $T_xN$ is top-linear, the averaging is also top-linear.

\begin{Lemma}
 $\ker (\ave_x)$ and $(T_xN)^{\Gamma}$ are closed subspaces of $T_xN$.
 Furthermore, a~sequence
 \[0 \to \ker(\ave_x) \hookrightarrow T_xN \twoheadrightarrow (T_xN)^{\Gamma} \to 0\]
 is split.
\end{Lemma}

\begin{proof}
 Since $(T_xN)^{\Gamma}$ is an intersection of closed subspaces \smash{$(T_xN)^{\Gamma} = \bigcap_{\varphi \in \Gamma} \ker(\varphi - 1)$},
 the space $(T_xN)^{\Gamma}$ is closed.
 Because $\ave_x$ is top-linear, the kernel $\ker(\ave_x)$ is a~closed subspace.
 For any $X \in T_xN$,
 \begin{align*}
 (\ave_x)^2(X) &= \ave_x\bigg(\frac{1}{|\Gamma|}\sum_{\varphi \in \Gamma}\varphi\cdot X \bigg)
 = \frac{1}{|\Gamma|}\sum_{\varphi \in \Gamma} \ave_x(\varphi \cdot X)
 = \frac{1}{|\Gamma|} \sum_{\varphi \in \Gamma} \varphi \cdot X
 = \ave_x(X).
 \end{align*}
 Therefore, we have $T_xN = \ker(\ave_x) \oplus (T_xN)^{\Gamma}$.
\end{proof}

Suppose that $N^{\Gamma}$ is connected. Fix an element $x \in N^{\Gamma}$.
We show that there exists a~good chart around $x \in N$.
By Bochner's linearization theorem for compact Lie group actions on Banach manifolds~\cite[Proposition~3.1.4]{diez2019normal}, there exists a~$\Gamma$-invariant open neighborhood $U$ of~${x \in N^{\Gamma}}$
and a~$\Gamma$-equivariant diffeomorphism $\chi \colon U \stackrel{\sim}{\rightarrow} V \subset T_xN $ which satisfies ${\rm d}_x\chi =1_{T_xN}$.
Since $T_xN = \ker(\ave_x) \oplus (T_xN)^{\Gamma}$, the above neighborhood $U$ can be written as $ U= \chi^{-1}\bigl(V^{\prime}\times V^{\Gamma}\bigr)$
where $V^{\prime}$ \big(resp.\ $V^{\Gamma}$\big) is a~sufficiently small open neighborhood of $0$ in $\ker(\ave_x)$ \big(resp.\ $(T_xN)^{\Gamma}$\big).
Because $\chi$ is $\Gamma$-equivariant, we have $N^{\Gamma} \cap U \simeq V^{\Gamma} = V \cap (T_xN)^{\Gamma}$.
Therefore, each connected component of $N^{\Gamma}$ is a~Banach submanifold of $N$.


\begin{thebibliography}{99}
\footnotesize\itemsep=0pt

\bibitem{ABM}
Alekseev A., Bursztyn H., Meinrenken E., Pure spinors on {L}ie groups,
 \textit{Ast\'{e}risque} \textbf{327} (2009), 131--199.

\bibitem{AMM}
Alekseev A., Malkin A., Meinrenken E., Lie group valued moment maps,
 \href{https://doi.org/10.4310/jdg/1214460860}{\textit{J.~Differential Geom.}}
 \textbf{48} (1998), 445--495,
 \href{https://arxiv.org/abs/dg-ga/9707021}{arXiv:dg-ga/9707021}.

\bibitem{Atiyah-Bott}
Atiyah M.F., Bott R., The {Y}ang--{M}ills equations over {R}iemann surfaces,
 \href{https://doi.org/10.1098/rsta.1983.0017}{\textit{Philos. Trans. Roy.
 Soc. London Ser.~A}} \textbf{308} (1983), 523--615.

\bibitem{Boalch_alg_qHam}
Boalch P., Quasi-{H}amiltonian geometry of meromorphic connections,
 \href{https://doi.org/10.1215/S0012-7094-07-13924-3}{\textit{Duke Math.~J.}}
 \textbf{139} (2007), 369--405,
 \href{https://arxiv.org/abs/math/0203161}{arXiv:math/0203161}.

\bibitem{B-Y2015}
Boalch P., Yamakawa D., Twisted wild character varieties,
 \href{https://arxiv.org/abs/1512.08091}{arXiv:1512.08091}.

\bibitem{BHCM}
Bursztyn H., Crainic M., Dirac geometry, quasi-{P}oisson actions and
 {$D/G$}-valued moment maps,
 \href{https://doi.org/10.4310/jdg/1251122545}{\textit{J.~Differential Geom.}}
 \textbf{82} (2009), 501--566,
 \href{https://arxiv.org/abs/0710.0639}{arXiv:0710.0639}.

\bibitem{diez2019normal}
Diez T., Normal form of equivariant maps and singular symplectic reduction in
 infinite dimensions with applications to gauge field theory, Ph.D.~Thesis,
 {U}niversit\"at {L}eipzig, 2019,
 \href{https://arxiv.org/abs/1909.00744}{arXiv:1909.00744}.

\bibitem{Donaldson1992}
Donaldson S.K., Boundary value problems for {Y}ang--{M}ills fields,
 \href{https://doi.org/10.1016/0393-0440(92)90044-2}{\textit{J.~Geom. Phys.}}
 \textbf{8} (1992), 89--122.

\bibitem{F-R}
Fock V.V., Rosly A.A., Poisson structure on moduli of flat connections on
 {R}iemann surfaces and the {$r$}-matrix, in Moscow {S}eminar in
 {M}athematical {P}hysics, \textit{Amer. Math. Soc. Transl. Ser.~2}, Vol.~191,
 \href{https://doi.org/10.1090/trans2/191/03}{American Mathematical Society},
 Providence, RI, 1999, 67--86.

\bibitem{Goldman}
Goldman W.M., The symplectic nature of fundamental groups of surfaces,
 \href{https://doi.org/10.1016/0001-8708(84)90040-9}{\textit{Adv. Math.}}
 \textbf{54} (1984), 200--225.

\bibitem{YealKarshon}
Karshon Y., An algebraic proof for the symplectic structure of moduli space,
 \href{https://doi.org/10.2307/2159424}{\textit{Proc. Amer. Math. Soc.}}
 \textbf{116} (1992), 591--605.

\bibitem{Knop}
Knop F., Classification of multiplicity free quasi-{H}amiltonian manifolds,
 \href{https://doi.org/10.4310/pamq.2024.v20.n1.a10}{\textit{Pure Appl.
 Math.~Q.}} \textbf{20} (2024), 471--523,
 \href{https://arxiv.org/abs/2210.07637}{arXiv:2210.07637}.

\bibitem{LYM_twisted_LG}
Loizides Y., Meinrenken E., Song Y., Spinor modules for {H}amiltonian loop
 group spaces,
 \href{https://doi.org/10.4310/JSG.2020.v18.n3.a10}{\textit{J.~Symplectic
 Geom.}} \textbf{18} (2020), 889--937,
 \href{https://arxiv.org/abs/1706.07493}{arXiv:1706.07493}.

\bibitem{mac1998categories}
Mac~Lane S., Categories for the working mathematician, 2nd ed., \textit{Grad. Texts in
 Math.}, Vol.~5,
 \href{https://doi.org/10.1007/978-1-4757-4721-8}{Springer}, New York, 1998.


\bibitem{Meinrenken2017}
Meinrenken E., Convexity for twisted conjugation,
 \href{https://doi.org/10.4310/MRL.2017.v24.n6.a12}{\textit{Math. Res. Lett.}}
 \textbf{24} (2017), 1797--1818,
 \href{https://arxiv.org/abs/1512.09000}{arXiv:1512.09000}.

\bibitem{Meinrenken2025}
Meinrenken E., Introduction to moduli spaces and {D}irac geometry,
 \href{https://arxiv.org/abs/2506.04150}{arXiv:2506.04150}.

\bibitem{Zerouali_Ph.D_Thesis}
Zerouali A.J., Twisted conjugation, quasi-{H}amiltonian geometry, and
 {D}uistermaat--{H}eckman measures, Ph.D.~Thesis, {U}niversity of Toronto,
 2019.

\bibitem{Zerouali}
Zerouali A.J., Twisted moduli spaces and {D}uistermaat--{H}eckman measures,
 \href{https://doi.org/10.1016/j.geomphys.2020.104042}{\textit{J.~Geom.
 Phys.}} \textbf{161} (2021), 104042, 27~pages,
 \href{https://arxiv.org/abs/2007.00103}{arXiv:2007.00103}.

\end{thebibliography} 
\end{document}
